\section{Discussion}
Network proximity is a central concept of network medicine, but its use for functional enrichment has remained ad hoc: proximity scores mix annotation overlap with network evidence, significance rests on a limited number of permutations and a normal approximation, and the interactome may itself encode the pathway annotations being tested. NetFunrankR v2 turns proximity-based enrichment into a statistically explicit procedure. The key observation is that, under the standard degree-preserving null model of network medicine, overlap, distance and linker statistics are sums of independent discrete variables, so their null distributions can be computed exactly. Conditioning on the membership pattern then separates the evidence into an overlap test and beyond-overlap tests that are independent under the null (Proposition~\ref{prop:indep}). This yields (i) a valid combined test, (ii) an attribution of each finding to direct annotation, network context, or both, and (iii) p-values that are exact far into the tail at a fraction of the cost of permutation testing.

Three results stand out from the benchmarks. First, calibration matters more than is commonly appreciated. Random gene lists that merely over-represent highly connected proteins---a property shared by most disease-gene and GWAS lists---made the hypergeometric ORA report FDR-significant pathways for almost every list of 100 genes or more, and NEA did so for every list regardless of composition. A degree-aware overlap test removes this artefact and, perhaps surprisingly, also improves the ranking of phenotype-relevant pathways. Second, network evidence is most valuable exactly where annotation-based enrichment is weakest: for incomplete annotation (simulation), for GWAS gene lists, and for short lists, where ORA recovered less than 10\% of the attainable COPD relevance from 6--10 genes while NetFunrankR recovered more than half. Third, among beyond-overlap statistics the number of linker genes was clearly superior to mean distance for genetic data, consistent with sparse signals in which only a minority of loci relate to a given pathway; the distance statistic, which underlies previous proximity analyses including NetFunrankR v1, was the weakest component.

Several limitations should be noted. The null model treats query genes as independent draws within degree bins; it does not model linkage disequilibrium between GWAS loci, gene length, or co-regulation of differentially expressed genes, which can induce dependence among query genes. Like ORA, NetFunrankR is a competitive test of a gene list and cannot exploit the continuous statistics or sample structure that self-contained expression methods use; on expression data it should therefore be viewed as complementary to methods such as CAMERA or PADOG rather than as a replacement. MalaCards relevance rankings, the gold standard of GSEABenchmarkeR, are partly derived from gene overlap and are incomplete, which favours overlap-based methods and limits the attainable precision of any method that reports pathways supported by network evidence only; the simulations and the temporal validation were designed to complement this. The choice of the linker statistic as the default beyond-overlap test was made on the GWAS benchmark, and we therefore report both variants throughout; the other benchmarks were not used for method selection. Finally, interactome coverage remains biased towards well-studied proteins. Removing the curated-database and text-mining channels of STRING eliminates direct circularity but not this study bias, which the degree-preserving null mitigates but cannot remove entirely.

Beyond pathway enrichment, the same construction applies whenever a gene list is compared with a gene module in a network, for example drug targets versus disease modules \citep{Guney2016} or disease--disease relationships \citep{Menche2015}, where exact, decomposed and calibrated tests could replace permutation $z$-scores.
